\scanpage{165}
$V$ un $\unsure{G^0|G^*}$-module \inserted{de poids $\unsure{p|\rho}$ (N.B. le \unsure{norme} est de poids $-2$)}, $V_{\mathbf C}=V\otimes_{\mathbf Q}\mathbf C$

Alors une place \deleted{réelle} $\xi$ permet de donner un sous-\inserted{$\unsure{\widetilde K|\overline K}$-}espace vectoriel $H_\xi(V)\subset V_{\mathbf C}$, filtré par des $H_\xi(\unsure{W|V})^{(p)}$, tel que \deleted{\unsure{ce}} ce soit une \uline{réduction des scalaires} à $K$, et que la filtration $H_\xi(V)_{\mathbf C}^{(p)}$ fasse de $V$ un \uuline{réseau de} \uline{Hodge}, $G^0$ étant \deleted{alors} \inserted{exactement} le \uline{groupe de Galois} de la structure de Hodge. Si la place est réelle, alors
\margin{\unsure{coh. de De Rham}}

\deleted{\unsure{On}} $V$ un $\unsure{G|C}$-module,

$\xi$ une place de $K$
\[
\overline V_\xi=V\otimes_{\mathbf Q}\overline K_\xi\supset H_{\unsure{\xi}}(V)^{(p)}
\]
\margin{\diagram{\begin{tikzcd}[column sep=small,row sep=small]
K_{\xi_1} \arrow[dr,no head] & K_{\xi_2} \arrow[d,no head] & K_{\xi_{\unsure{m|\infty}}} \arrow[dl,no head] \\
& K \arrow[d,no head] & \\
& \mathbf Q &
\end{tikzcd}}}
Les $H_\xi(V)^p\otimes_K K_\xi$ définissent sur $(V,V_K)$ \uline{une structure de Hodge}, \deleted{\uline{\unsure{le}}} \uline{groupe de Galois étant} \deleted{\unsure{l'image de}} \uline{\deleted{\unsure{tous}} $\operatorname{Im}G^0$}

Si $\xi$ est réelle, on a une involution dans $V$ provenant de l'élément \deleted{$\unsure{f_\xi}$} $\operatorname{frob}_\xi$ de $G(\mathbf Q)$ et on veut que \unsure{prolongeant} $\operatorname{frob}_\xi$ de façon \uline{$\unsure{K_\xi}$-linéaire} à $\overline V_\xi$, $\operatorname{frob}'_\xi$ soit l'identité sur $H_\xi(V)$ [et par suite $\operatorname{frob}_\xi$ est compatible avec la structure de Hodge $\ldots$]
\margin{164}
